\MulaiModul{09}{Pencetakan bilangan}

\subsection{Tujuan}

\begin{tujuan}
Setelah menyelesaikan modul ini, mahasiswa mampu:
\begin{enumerate}
    \item menjelaskan mengapa bilangan harus diubah menjadi deret
          digit sebelum dicetak;
    \item menggunakan \texttt{DIV} untuk pembagian byte dan word
          serta menjelaskan hasil bagi dan sisanya;
    \item mengosongkan \texttt{DX} sebelum pembagian word;
    \item mencetak satu bilangan tidak bertanda lebih dari satu digit
          dalam bentuk desimal;
    \item mencetak satu word sebagai empat digit heksadesimal;
    \item memakai prosedur untuk memecah rutin pencetakan.
\end{enumerate}
\end{tujuan}

\subsection{Dasar singkat}

\subsubsection*{Mengapa bilangan perlu diubah}

Layanan \texttt{AH=02h} hanya mencetak satu karakter. Isi register
berupa bilangan biner, bukan kode ASCII. Karena itu, bilangan harus
diubah lebih dahulu menjadi deret digit, lalu setiap digit diubah
menjadi kode ASCII. Pencetakan bilangan dikerjakan sendiri oleh
program; DOS tidak menyediakan layanan yang langsung mencetak
bilangan.

\subsubsection*{Instruksi DIV}

\texttt{DIV} membagi bilangan pada \texttt{AX} atau \texttt{DX:AX}
dengan operand pembagi. Hasil bagi dan sisa diletakkan pada register
yang berbeda menurut ukuran pembagi, seperti ditunjukkan pada
Tabel~\ref{tab:m09-div}.

\begin{table}[H]
    \centering
    \caption{Hasil \texttt{DIV} menurut ukuran pembagi.}
    \label{tab:m09-div}
    \small
    \begin{tabular}{@{}llll@{}}
        \toprule
        \textbf{Pembagi} & \textbf{Yang dibagi} &
        \textbf{Hasil bagi} & \textbf{Sisa} \\
        \midrule
        8 bit & \texttt{AX} & \texttt{AL} & \texttt{AH} \\
        16 bit & \texttt{DX:AX} & \texttt{AX} & \texttt{DX} \\
        \bottomrule
    \end{tabular}
\end{table}

Pada pembagian word, \texttt{DX:AX} dipakai sebagai bilangan yang
dibagi. Sebelum \texttt{DIV}, \texttt{DX} harus dikosongkan agar
bilangan yang dibagi hanya berisi nilai pada \texttt{AX}. Bila
\texttt{DX} tidak dikosongkan, sisa lama pada \texttt{DX} ikut
menjadi bagian bilangan dan hasil pembagian menjadi salah. Contoh:
dengan \texttt{AX=12345}, \texttt{BX=10}, dan \texttt{DX=0},
\texttt{DIV BX} menghasilkan hasil bagi \texttt{1234} pada
\texttt{AX} dan sisa \texttt{5} pada \texttt{DX}.

\subsubsection*{Konversi desimal}

Algoritma desimal membagi bilangan dengan 10 secara berulang. Setiap
sisa pembagian disimpan di stack sebagai satu digit, dan cacah digit
dicatat pada \texttt{CX}. Pembagian diulang sampai hasil bagi nol.
Setelah itu, digit diambil dari stack satu per satu, ditambah
\texttt{'0'} agar menjadi kode ASCII, lalu dicetak.

Sisa pembagian muncul dari digit paling kanan ke paling kiri,
sedangkan pencetakan harus dari kiri ke kanan. Sifat LIFO pada stack
membalik urutan itu sehingga digit tercetak dengan benar. Untuk
bilangan nol, pembagian tetap dijalankan sekali sehingga satu digit
\texttt{'0'} tersimpan dan dicetak.

\subsubsection*{Konversi heksadesimal}

Heksadesimal memakai basis 16 dan empat bit per digit. Prinsip
pembagian yang sama dipakai dengan pembagi 16. Setiap sisa bernilai
0 sampai 15. Sisa 0--9 diubah dengan menambahkan \texttt{'0'},
sedangkan sisa 10--15 diubah dengan menambahkan 7 lebih dahulu, baru
kemudian \texttt{'0'}. Penambahan 7 itu memindahkan 10--15 ke posisi
\texttt{'A'}--\texttt{'F'}. Tabel~\ref{tab:m09-digit} merangkumnya.

\begin{table}[H]
    \centering
    \caption{Konversi sisa pembagian menjadi kode ASCII.}
    \label{tab:m09-digit}
    \small
    \begin{tabular}{@{}lll@{}}
        \toprule
        \textbf{Nilai sisa} & \textbf{Perlakuan} & \textbf{Contoh} \\
        \midrule
        0--9 & tambah \texttt{30h} & 5 menjadi \texttt{35h} = \texttt{'5'} \\
        10--15 & tambah \texttt{37h} & 10 menjadi \texttt{41h} = \texttt{'A'} \\
        \bottomrule
    \end{tabular}
\end{table}

Rutin heksadesimal dijalankan empat kali sehingga satu word selalu
dicetak sebagai empat digit, termasuk angka nol di depan. Prosedur
pencetakan menyimpan dan memulihkan register yang dipakainya dengan
\texttt{PUSH} dan \texttt{POP} agar nilai pemanggil tidak berubah.

\begin{rujukan}
Rujukan singkat: Gunawan, Bagian 14.1.1 dan 14.1.2; S'to, Bab XX;
buku ajar, berkas \texttt{cetak-bilangan.asm} dan \texttt{chapters/15}.
\end{rujukan}

\subsection{Alat dan bahan}

\begin{itemize}
    \item Komputer dengan emu8086 terpasang.
    \item DOSBox dengan TASM, TLINK, dan DEBUG.
    \item Berkas \texttt{code/m09-cetak-desimal.asm} dan
          \texttt{code/m09-cetak-hex.asm}.
    \item Lembar jurnal praktikum.
\end{itemize}

\subsection{Contoh program}

Program pertama mencetak \texttt{12345} dalam bentuk desimal.
Prosedur \texttt{CETAK\_DESIMAL} membagi dengan 10, menyimpan sisa di
stack, lalu mencetak digit dari kiri ke kanan.

\lstinputlisting[
    caption={Mencetak bilangan desimal lebih dari satu digit.},
    label={lst:m09-desimal}
]{code/m09-cetak-desimal.asm}

Program kedua mencetak word yang sama sebagai heksadesimal. Karena
\texttt{12345 = 3039h}, keluarannya adalah \texttt{3039}. Prosedur
\texttt{CETAK\_HEX} membagi dengan 16 sebanyak empat kali.

\lstinputlisting[
    caption={Mencetak satu word sebagai empat digit heksadesimal.},
    label={lst:m09-hex}
]{code/m09-cetak-hex.asm}

\subsection{Langkah kerja}

\begin{enumerate}
    \item Siapkan folder kerja, lalu salin kedua berkas \texttt{.asm}
          ke dalamnya.
    \item Rakit dan tautkan program pertama. Jalankan kedua perintah
          berikut pada baris terpisah:
          \begin{perintah}
tasm m09-cetak-desimal.asm\\
tlink m09-cetak-desimal.obj
\end{perintah}
          Jalankan hasilnya dan catat bilangan yang tercetak.
    \item Ubah \texttt{NILAI} menjadi 0, rakit ulang, lalu jalankan.
          Ulangi dengan \texttt{NILAI} bernilai 65535. Catat kedua
          keluarannya.
    \item Jalankan program pertama dengan \emph{single step} pada
          emu8086. Pada setiap \texttt{DIV}, catat nilai \texttt{AX}
          dan \texttt{DX} sebelum serta sesudah pembagian. Pastikan
          \texttt{DX} dikosongkan lebih dahulu.
    \item Rakit dan tautkan program kedua dengan perintah yang sama,
          lalu jalankan. Catat keluaran empat digitnya.
    \item Ubah \texttt{NILAI} menjadi \texttt{000Fh}, lalu menjadi
          \texttt{0FFFFh}. Rakit ulang dan jalankan kedua versi.
          Jelaskan bagaimana sisa 10--15 diubah menjadi
          \texttt{'A'}--\texttt{'F'}.
\end{enumerate}

\subsection{Tugas}

\begin{enumerate}
    \item Jelaskan hasil \texttt{DIV} untuk pembagi 8 bit dan pembagi
          16 bit. Sebutkan register hasil bagi dan sisa pada
          masing-masing kasus.
    \item Jelaskan mengapa \texttt{DX} harus dikosongkan sebelum
          \texttt{DIV} word dan apa akibatnya bila tidak dikosongkan.
    \item Jelaskan algoritma konversi desimal dan peran stack di
          dalamnya. Telusuri sisa pembagian untuk bilangan 243.
    \item Jelaskan cara mengubah sisa 10--15 menjadi karakter
          \texttt{'A'}--\texttt{'F'}.
    \item Tentukan deret digit desimal dan heksadesimal untuk bilangan
          12345, lalu bandingkan hasilnya dengan keluaran kedua
          program.
\end{enumerate}

\subsection{Hasil yang dikumpulkan}

\begin{enumerate}
    \item Kedua berkas \texttt{.asm} beserta versi modifikasinya.
    \item Tangkapan layar keluaran kedua program untuk beberapa nilai
          \texttt{NILAI}.
    \item Tabel jejak \texttt{DIV} berisi \texttt{AX} dan \texttt{DX}
          sebelum serta sesudah pembagian.
    \item Lembar jurnal berisi jawaban langkah kerja dan tugas.
\end{enumerate}
